% Fuente bilingüe común: las fórmulas y pruebas estructurales son idénticas.
\clearpage
\section{\HMTLang{Transporte simpléctico de la elipse de acción}{Symplectic transport of the action ellipse}}
\label{delta22:ii:transporte}
\HMTLang{La elipse previamente construida permite distinguir la medida areal de acción, la anisotropía angular y la tasa de evolución. Esta distinción determina un transporte exacto entre realizaciones elípticas y fija el término de conexión de su expresión hamiltoniana. El antecedente es el par angular publicado por APP--TRIT--TPK, después de la generación regional y de la selección conjunta de \(K\); las operaciones siguientes actúan sobre esas salidas y conservan sus etiquetas de orientación.}{The ellipse constructed above distinguishes the areal measure of action, angular anisotropy, and the rate of evolution. This distinction determines an exact transport between elliptical realizations and fixes the connection term in its Hamiltonian expression. Its antecedent is the angular pair expressed by APP--TRIT--TPK after regional generation and the joint selection of \(K\); the following operations act on those outputs and retain their orientation labels.}

\subsection{\HMTLang{Forma cuadrática y realización inductivo--capacitiva}{Quadratic form and inductive--capacitive realization}}
\HMTLang{Escribimos \(\zeta=\eta_{\rm el}=\operatorname{artanh}(C^*/A)\), con ambas coordenadas angulares en la misma unidad y \(|C^*/A|<1\). El símbolo \(\eta_{\rm ret}\) queda reservado al coeficiente de acción. Para la sección positiva \(\hbar=\eta_{\rm ret}s_0\), \(s_0=10^{-34}\mathcal U_S\), las coordenadas equilibradas de dimensión raíz de acción son}{Write \(\zeta=\eta_{\rm el}=\operatorname{artanh}(C^*/A)\), with both angular coordinates in the same unit and \(|C^*/A|<1\). The symbol \(\eta_{\rm ret}\) is reserved for the action coefficient. For the positive section \(\hbar=\eta_{\rm ret}s_0\), \(s_0=10^{-34}\mathcal U_S\), the balanced coordinates, with the dimension of the square root of action, are}
\[
 Q=\sqrt{2\hbar}\,e^{\zeta/2}\cos\theta,
 \qquad P=\sqrt{2\hbar}\,e^{-\zeta/2}\sin\theta.
\]
\HMTLang{El producto de los semiejes es \(2\hbar\), el área total es \(2\pi\hbar=h\) y el área paramétrica acumulada es \(\hbar\theta\). La fase \(\theta\) es el parámetro de esta representación elíptica. Conservando el transductor \(Z_*>0\) y la frecuencia \(\omega>0\) de la realización anterior, se tiene}{The product of the semiaxes is \(2\hbar\), the total area is \(2\pi\hbar=h\), and the accumulated parametric area is \(\hbar\theta\). The phase \(\theta\) is the parameter of this elliptical representation. Retaining the transducer \(Z_*>0\) and the frequency \(\omega>0\) of the preceding realization gives}
\[
 L_\zeta=\frac{Z_*}{\omega}e^{-\zeta},\qquad
 C_\zeta=\frac{e^\zeta}{Z_*\omega},\qquad
 Q=\sqrt{Z_*}\,q,\quad P=\frac{\Phi}{\sqrt{Z_*}}.
\]
\begin{proposition}[\HMTLang{Composición energética de la elipse}{Energy composition of the ellipse}]
\[
 H_0=\frac{\omega}{2}(e^{-\zeta}Q^2+e^\zeta P^2),\quad
 E_{\rm el}=\hbar\omega\cos^2\theta,\quad
 E_{\rm mag}=\hbar\omega\sin^2\theta.
\]
\HMTLang{Sobre el contorno, \(H_0=E_{\rm el}+E_{\rm mag}=\hbar\omega\).}{On the boundary, \(H_0=E_{\rm el}+E_{\rm mag}=\hbar\omega\).}
\end{proposition}
\begin{proof}
\HMTLang{La sustitución de \(q=Q/\sqrt{Z_*}\), \(\Phi=\sqrt{Z_*}P\) en \(q^2/(2C_\zeta)+\Phi^2/(2L_\zeta)\) produce \(H_0\). Sustituir después la parametrización elíptica da las dos contribuciones y su suma. La acción y la forma preceden al circuito que las realiza.}{Substitution of \(q=Q/\sqrt{Z_*}\), \(\Phi=\sqrt{Z_*}P\) into \(q^2/(2C_\zeta)+\Phi^2/(2L_\zeta)\) yields \(H_0\). Substituting the elliptical parametrization then gives the two contributions and their sum. Action and shape precede the circuit realizing them.}
\end{proof}
\HMTLang{La impedancia de esta realización satisface \(\sqrt{L_\zeta/C_\zeta}/Z_*=e^{-\zeta}\). La impedancia constitutiva del vacío utiliza el cociente de sus lectores \(90/120\); cada impedancia conserva su definición y su transducción.}{The impedance of this realization satisfies \(\sqrt{L_\zeta/C_\zeta}/Z_*=e^{-\zeta}\). The constitutive impedance of the vacuum uses the ratio of its \(90/120\) readers; each impedance retains its definition and transduction.}

\subsection{\HMTLang{Invariante de transporte y conexión hamiltoniana}{Transport invariant and Hamiltonian connection}}
\HMTLang{Sean}{Let}
\[
 S_\zeta=\operatorname{diag}(e^{\zeta/2},e^{-\zeta/2}),\quad
 g_\zeta=S_\zeta^{-\mathsf T}S_\zeta^{-1},\quad
 J_0=\begin{pmatrix}0&-1\\1&0\end{pmatrix},\quad
 R(-\theta)=e^{-\theta J_0}.
\]
\begin{theorem}[\HMTLang{Conservación exacta de la acción cuadrática}{Exact conservation of quadratic action}]
\HMTLang{Para una sucesión de formas \(\zeta_n\) y fases \(\theta_n\), el transporte}{For a sequence of shapes \(\zeta_n\) and phases \(\theta_n\), the transport}
\[
 X_{n+1}=T_nX_n,\qquad
 T_n=S_{\zeta_{n+1}}R(-\theta_n)S_{\zeta_n}^{-1}
\]
\HMTLang{tiene determinante uno y conserva}{has determinant one and preserves}
\[
 I_n=\tfrac12X_n^{\mathsf T}g_{\zeta_n}X_n,
 \qquad T_n^{\mathsf T}g_{\zeta_{n+1}}T_n=g_{\zeta_n}.
\]
\end{theorem}
\begin{proof}
\HMTLang{Los determinantes de \(S_\zeta\) y \(R\) son uno. En el producto cuadrático, los factores \(S_{\zeta_{n+1}}\) cancelan sus inversos; \(R^{\mathsf T}R=I\) deja \(S_{\zeta_n}^{-\mathsf T}S_{\zeta_n}^{-1}\). Esta identidad prueba la conservación para cada paso y, por composición, para cualquier longitud. Una sucesión publicada por una ruta TPK se transporta con esa ruta y sus etiquetas.}{The determinants of \(S_\zeta\) and \(R\) are one. In the quadratic product, the factors \(S_{\zeta_{n+1}}\) cancel their inverses; \(R^{\mathsf T}R=I\) leaves \(S_{\zeta_n}^{-\mathsf T}S_{\zeta_n}^{-1}\). This identity proves conservation at every step and, by composition, at any length. A sequence expressed by a TPK route is transported together with that route and its labels.}
\end{proof}
\begin{theorem}[\HMTLang{Generador diferenciable y término mixto}{Differentiable generator and mixed term}]
\HMTLang{Si \(\zeta\) es de clase \(C^1\), \(\omega>0\) es continua y \(X=S_\zeta Y\), \(\dot Y=-\omega J_0Y\), entonces}{If \(\zeta\) is \(C^1\), \(\omega>0\) is continuous, and \(X=S_\zeta Y\), \(\dot Y=-\omega J_0Y\), then}
\[
 \dot Q=\tfrac12\dot\zeta Q+\omega e^\zeta P,
 \qquad \dot P=-\omega e^{-\zeta}Q-\tfrac12\dot\zeta P.
\]
\HMTLang{Con la convención \(\dot Q=\partial_PH\), \(\dot P=-\partial_QH\), su Hamiltoniano, salvo un sumando dependiente del tiempo, es}{With the convention \(\dot Q=\partial_PH\), \(\dot P=-\partial_QH\), its Hamiltonian, up to an additive function of time, is}
\begin{equation}
 H_{\rm tot}=\frac\omega2(e^{-\zeta}Q^2+e^\zeta P^2)
             +\frac{\dot\zeta}{2}QP.
 \label{delta22:ii:hamiltoniano}
\end{equation}
\HMTLang{El invariante \(I=(e^{-\zeta}Q^2+e^\zeta P^2)/2\) se conserva exactamente.}{The invariant \(I=(e^{-\zeta}Q^2+e^\zeta P^2)/2\) is conserved exactly.}
\end{theorem}
\begin{proof}
\HMTLang{Derivar \(X=S_\zeta Y\) produce \(\dot S_\zeta S_\zeta^{-1}=\operatorname{diag}(\dot\zeta/2,-\dot\zeta/2)\) y el término rotacional conjugado. Las derivadas parciales de \eqref{delta22:ii:hamiltoniano} reproducen las dos ecuaciones. En \(\dot I\), las contribuciones explícitas \(-\dot\zeta e^{-\zeta}Q^2/2+\dot\zeta e^\zeta P^2/2\) cancelan las procedentes de la diagonal del generador; los términos cruzados rotacionales se anulan entre sí.}{Differentiating \(X=S_\zeta Y\) yields \(\dot S_\zeta S_\zeta^{-1}=\operatorname{diag}(\dot\zeta/2,-\dot\zeta/2)\) and the conjugated rotational term. The partial derivatives of \eqref{delta22:ii:hamiltoniano} reproduce both equations. In \(\dot I\), the explicit contributions \(-\dot\zeta e^{-\zeta}Q^2/2+\dot\zeta e^\zeta P^2/2\) cancel those from the diagonal of the generator; the rotational cross terms cancel each other.}
\end{proof}
\HMTLang{En la parametrización elíptica, \(QP=I\sin2\theta\) y \(\dot\theta=-\omega\). Por tanto \(H_{\rm tot}=I\omega+(I\dot\zeta/2)\sin2\theta\). La matriz cuadrática tiene determinante \(\omega^2-\dot\zeta^2/4\): es definida positiva exactamente para \(|\dot\zeta|<2\omega\), y semidefinida positiva en igualdad. La conservación de \(I\) y la dependencia temporal de \(H_{\rm tot}\) expresan propiedades distintas. Como control algebraico, omitir el término mixto produce \(\dot I=\dot\zeta(e^\zeta P^2-e^{-\zeta}Q^2)/2\). En la representación cuántica correspondiente, la expresión simétrica del término mixto es \(\dot\zeta(\widehat Q\widehat P+\widehat P\widehat Q)/4\).}{In the elliptical parametrization, \(QP=I\sin2\theta\) and \(\dot\theta=-\omega\). Thus \(H_{\rm tot}=I\omega+(I\dot\zeta/2)\sin2\theta\). The quadratic matrix has determinant \(\omega^2-\dot\zeta^2/4\): it is positive definite precisely when \(|\dot\zeta|<2\omega\), and positive semidefinite at equality. Conservation of \(I\) and time dependence of \(H_{\rm tot}\) express different properties. As an algebraic control, omitting the mixed term gives \(\dot I=\dot\zeta(e^\zeta P^2-e^{-\zeta}Q^2)/2\). In the corresponding quantum representation, the symmetric mixed term is \(\dot\zeta(\widehat Q\widehat P+\widehat P\widehat Q)/4\).}

\subsection{\HMTLang{Compatibilidad genealógica y transporte conformemente simpléctico}{Genealogical compatibility and conformally symplectic transport}}
\HMTLang{En la sección generada,}{On the generated section,}
\[
 \zeta=\operatorname{artanh}\frac{2(\eta_{\rm ret}+\alpha)}{1000\alpha}.
\]
\HMTLang{Fijar \(\alpha\) y \(\eta_{\rm ret}\) fija la forma. La versión variable del teorema describe el transporte entre formas admitidas; su realización dinámica conserva la actualización que determina esas formas. Para secciones positivas \(\hbar_n\) en una misma carta dimensional, sea \(\rho_n=\hbar_{n+1}/\hbar_n\). Entonces}{Fixing \(\alpha\) and \(\eta_{\rm ret}\) fixes the shape. The variable version of the theorem describes transport between admissible shapes; its dynamical realization retains the update determining those shapes. For positive sections \(\hbar_n\) in the same dimensional chart, let \(\rho_n=\hbar_{n+1}/\hbar_n\). Then}
\[
 \widetilde T_n=\sqrt{\rho_n}\,T_n,\quad
 \det\widetilde T_n=\rho_n,\quad
 \widetilde T_n^{\mathsf T}g_{\zeta_{n+1}}\widetilde T_n
 =\rho_ng_{\zeta_n},\quad
 \frac{I_{n+1}}{\hbar_{n+1}}=\frac{I_n}{\hbar_n}.
\]
\HMTLang{Estas igualdades se deducen multiplicando el transporte anterior por \(\sqrt{\rho_n}\). La forma simpléctica y la acción normalizada se transportan conjuntamente. El límite diferenciable tiene traza \(\dot\hbar/\hbar\); en una forma simpléctica fija, un generador hamiltoniano lineal tiene traza cero. Por ello, la sección variable exige conservar la forma transportada o los grados de libertad que intercambian acción. Si la variación es un cambio de unidades, se transportan primero las bases dimensionales.}{These equalities follow by multiplying the preceding transport by \(\sqrt{\rho_n}\). The symplectic form and normalized action are transported together. The differentiable limit has trace \(\dot\hbar/\hbar\); for a fixed symplectic form, a linear Hamiltonian generator has zero trace. A variable section therefore requires retaining the transported form or the degrees of freedom exchanging action. If the variation is a change of units, the dimensional bases are transported first.}

\HMTLang{La distinción operacional es precisa. Bajo un cambio pasivo \(X=S_\zeta Y\), también el lector satisface \(L_X(S_\zeta Y)=L_Y(Y)\); el término mixto representa la conexión de coordenadas. Una deformación física se especifica mediante la evolución relativa entre sistema y referencia, obtenida de sus actualizaciones y leída con el mismo dispositivo. Una atribución energética incorpora el campo, el registro y el controlador. Así quedan separadas la conservación demostrada de la acción cuadrática, la realización hamiltoniana y las condiciones de una respuesta física contrastable.}{The operational distinction is precise. Under a passive change \(X=S_\zeta Y\), the reader also satisfies \(L_X(S_\zeta Y)=L_Y(Y)\); the mixed term represents the coordinate connection. A physical deformation is specified by relative evolution between system and reference, obtained from their updates and read with the same device. An energy attribution includes the field, the record, and the controller. This separates the proved conservation of quadratic action, its Hamiltonian realization, and the conditions for a testable physical response.}
